\chapter{Вложенные классы, исключения, обобщения и отладка}
\label{ch:04}

Освоив классы, интерфейсы и наследование, вы умеете описывать предметную
область. Эта глава добавляет четыре механизма, без которых не обходится ни один
настоящий проект: вложенные классы, чтобы вспомогательные сущности не
расползались по пакету; обобщения~--- способ писать код, работающий с любым
типом и при этом типобезопасный; исключения, которыми программа сообщает о
проблеме; и отладчик, нужный ровно тогда, когда программа не падает, а тихо
считает неправильно.

\section{Вложенные классы}

Представьте, что вы проектируете класс \texttt{Order} (заказ) с внутренним
понятием \texttt{OrderItem}~--- позицией заказа, которая отдельно от заказа не
имеет смысла. Выносить её в отдельный файл значит засорять пакет классом,
которым больше никто не пользуется. Правильнее объявить \texttt{OrderItem}
внутри \texttt{Order}: связанные понятия лежат рядом, а вложенный класс
получает доступ к закрытым полям внешнего.

Вложенные классы бывают двух видов: \term{нестатический внутренний класс}
(\eng{inner class}), объявленный без слова \texttt{static}, и \term{статический
вложенный класс} (\eng{static nested class}), объявленный с ним. Разница одна,
но принципиальная: внутренний класс хранит невидимую ссылку на объект внешнего
класса, а статический вложенный~--- нет.

\subsection{Нестатический внутренний класс}

Внутренний класс~--- как комната в доме: она не существует без дома, зато
пользуется всеми его ресурсами, то есть закрытыми полями внешнего класса, и
построить её раньше дома нельзя. Отсюда непривычный синтаксис создания~---
\texttt{outer.new Inner()}. В листинге~\ref{lst:04-1} собраны обе особенности;
обратите внимание на метод \texttt{show}, где три переменные с именем
\texttt{x} живут на трёх разных уровнях.

\begin{listing}[H]
\begin{minted}{java}
class Outer {
    private int x = 10;

    // Внутренний класс: ключевого слова static нет
    class Inner {
        private int x = 20;

        void show() {
            int x = 30;
            System.out.println(x);            // 30 - локальная
            System.out.println(this.x);       // 20 - поле Inner
            System.out.println(Outer.this.x); // 10 - поле Outer
        }
    }
}

// Сначала объект внешнего класса, только потом внутреннего
Outer outer = new Outer();
Outer.Inner inner = outer.new Inner();
\end{minted}
\caption{Внутренний класс и разрешение имён}
\label{lst:04-1}
\end{listing}

Конструкция \texttt{Outer.this} и есть та самая невидимая ссылка: компилятор
кладёт её в каждый объект \texttt{Inner} и позволяет обратиться к ней явно,
когда имена совпадают. Если имена различаются, поле внешнего класса видно
напрямую.

\begin{vrezka}
\textbf{Важно.} До Java~16 внутренний класс не мог объявлять статические члены,
кроме констант \texttt{static final} времени компиляции. Начиная с Java~16 это
ограничение снято.
\end{vrezka}

\subsection{Статический вложенный класс}

Статический вложенный класс~--- отдельный офис внутри здания компании: он в том
же здании, но работает самостоятельно и ключей от чужих кабинетов не имеет.
Доступ у него есть только к статическим членам внешнего класса, зато создавать
его можно без объекта внешнего~--- \texttt{new Outer.Nested()}.

Самое частое применение~--- паттерн \term{строитель} (\eng{Builder}). Он решает
проблему конструктора с длинным списком параметров, в котором легко перепутать
местами два соседних аргумента одного типа: объект собирают цепочкой вызовов с
говорящими именами (листинг~\ref{lst:04-2}). Каждый метод строителя возвращает
\texttt{this}, поэтому вызовы сцепляются, а конструктор класса закрыт, так что
недособранный объект получить невозможно.

\begin{listing}[H]
\begin{minted}{java}
public class Person {
    private final String firstName;
    private final int age;

    // Конструктор закрыт: объект собирается только строителем
    private Person(Builder b) {
        this.firstName = b.firstName;
        this.age = b.age;
    }

    public static class Builder {
        private String firstName;
        private int age;

        // Каждый метод возвращает сам строитель
        public Builder firstName(String v) {
            firstName = v; return this;
        }

        public Builder age(int v) {
            age = v; return this;
        }

        public Person build() { return new Person(this); }
    }
}

Person p = new Person.Builder().firstName("Иван").age(25).build();
\end{minted}
\caption{Статический вложенный класс: паттерн Builder}
\label{lst:04-2}
\end{listing}

Почему строитель именно статический? Потому что он создаёт объект внешнего
класса: требовать готовый \texttt{Person} ради того, чтобы собрать
\texttt{Person}, было бы замкнутым кругом. По той же причине статическим сделан
класс \texttt{Entry} внутри \texttt{HashMap}, а итератор внутри списка,
наоборот, внутренний: он обходит конкретный список и без него бессмыслен.
Различия видов сведены в табл.~\ref{tab:04-1}.

\begin{table}[htbp]
\caption{Внутренний и статический вложенный классы}
\label{tab:04-1}
\centering
\begin{tabular}{L{4.5cm}L{5.1cm}L{5.1cm}}
\toprule
Критерий & Внутренний (без \texttt{static}) & Статический вложенный \\
\midrule
Связь с объектом внешнего & есть неявная ссылка & нет \\
\addlinespace
Создание объекта & \texttt{outer.new Inner()} & \texttt{new Outer.Nested()} \\
\addlinespace
Доступ к нестатическим полям & да & нет \\
\addlinespace
Типичный пример & итератор коллекции & \texttt{Builder}, \texttt{Entry} \\
\bottomrule
\end{tabular}
\end{table}

\subsection{Вложенные интерфейсы}

Интерфейс тоже можно объявить внутри класса или внутри другого интерфейса, если
контракт логически принадлежит внешнему типу. Вложенный интерфейс всегда неявно
\texttt{static}~--- состояния у интерфейса нет, значит и привязывать его к
объекту незачем; внутри класса он может иметь любую видимость, а внутри
интерфейса всегда неявно \texttt{public static}. Реализовать его вправе кто
угодно, обратившись к нему через имя внешнего типа (листинг~\ref{lst:04-3}).

\begin{listing}[H]
\begin{minted}{java}
class Machine {
    // Неявно static, видимость задаём сами
    public interface PowerSwitch {
        void turnOn();
        void turnOff();
    }

    public static class Engine implements PowerSwitch {
        public void turnOn()  { System.out.println("Пуск"); }
        public void turnOff() { }
    }
}

// Посторонний класс тоже вправе реализовать контракт
class Lamp implements Machine.PowerSwitch {
    public void turnOn()  { }
    public void turnOff() { }
}
\end{minted}
\caption{Интерфейс, объявленный внутри класса}
\label{lst:04-3}
\end{listing}

\section{Обобщения}

До версии Java~5 коллекции хранили значения типа \texttt{Object}: положить в
список можно было что угодно, а при извлечении требовалось приведение типа~---
и если программист ошибался, компилятор молчал, а программа падала при запуске.
Сравните два фрагмента в листинге~\ref{lst:04-4}.

\begin{listing}[H]
\begin{minted}{java}
// Было: список хранит Object
List list = new ArrayList();
list.add(42);                     // никто не возражает
String s = (String) list.get(0);  // ClassCastException при запуске

// Стало: тип элементов известен компилятору
List<String> safe = new ArrayList<>();
safe.add("Строка");
// safe.add(42);                  // ошибка КОМПИЛЯЦИИ
String t = safe.get(0);           // приведение не нужно
\end{minted}
\caption{Коллекция без обобщений и с обобщениями}
\label{lst:04-4}
\end{listing}

\term{Обобщения} (\eng{generics}) дают три вещи: типобезопасность на этапе
компиляции, отсутствие ручных приведений и возможность написать один класс,
работающий с любым типом. Ошибка, пойманная компилятором, стоит дёшево; она же,
найденная пользователем в рабочей системе,~--- дорого.

\subsection{Параметры типа, ограничения и обобщённые методы}

Обобщённый класс объявляют, дописав к имени параметр типа в угловых скобках;
конкретное значение подставляется при создании объекта. Имена параметров
подчиняются соглашению: \texttt{T}~--- тип общего назначения (\eng{type}),
\texttt{E}~--- элемент коллекции, \texttt{K} и \texttt{V}~--- ключ и значение
отображения.

Иногда «любой тип» слишком широк: чтобы вызвать у значения метод
\texttt{doubleValue}, компилятор должен знать, что это число,~--- тогда
параметр ограничивают сверху словом \texttt{extends}. Обобщённым может быть и
отдельный метод: параметр типа объявляют перед возвращаемым значением. Все три
конструкции показаны в листинге~\ref{lst:04-5}; обратите внимание: у метода
\texttt{findMax} объявление \texttt{<T extends Comparable<T>>} стоит перед типом
результата, а не после имени метода.

\begin{listing}[H]
\begin{minted}{java}
// Обобщённый класс: T подставляется при создании объекта
class Box<T> {
    private T value;
    public void set(T value) { this.value = value; }
    public T get()           { return value; }
}

// Ограничение сверху: T - это Number или его наследник
class NumberBox<T extends Number> {
    private final T number;
    public NumberBox(T number) { this.number = number; }

    // Метод Number доступен, потому что граница объявлена
    public double asDouble() { return number.doubleValue(); }
}

// Обобщённый метод: <T> стоит перед типом результата
public static <T extends Comparable<T>> T findMax(T[] arr) {
    T max = arr[0];
    for (T e : arr) {
        if (e.compareTo(max) > 0) max = e;
    }
    return max;
}
\end{minted}
\caption{Обобщённый класс, ограничение типа и обобщённый метод}
\label{lst:04-5}
\end{listing}

Границ может быть несколько: запись
\texttt{<T extends Comparable<T> \& Cloneable>} требует выполнения обоих
контрактов сразу. Разделитель здесь~--- амперсанд, а не запятая: запятая
отделяет один параметр типа от другого.

\subsection{Стирание типов}

Ключевое свойство обобщений в Java~--- \term{стирание типов} (\eng{type
erasure}). Параметры типа существуют только на этапе компиляции: компилятор
проверяет по ним корректность кода, а затем удаляет их из байт-кода, заменяя на
\texttt{Object} или на объявленную верхнюю границу~--- ради совместимости с
кодом, написанным до Java~5.

Отсюда следствие, которое сначала удивляет: \texttt{List<String>} и
\texttt{List<Integer>} во время работы программы~--- один и тот же класс
\texttt{ArrayList}, и сравнение \texttt{strings.getClass() == ints.getClass()}
возвращает \texttt{true}. Из того же стирания вытекают запреты: внутри
обобщённого класса нельзя написать ни \texttt{new T()}, ни \texttt{new T[10]},
ни \texttt{obj instanceof T}, ни \texttt{T.class}~--- всё это требует знать тип
во время выполнения. Зато объявить поле типа \texttt{T}, параметр метода или
коллекцию \texttt{List<T>} можно: достаточно проверки при компиляции. Обходной
путь~--- хранить элементы в \texttt{Object[]} и приводить их при чтении:
\texttt{(T) array[i]}; компилятор выдаст предупреждение \eng{unchecked},
которое гасят аннотацией \texttt{@SuppressWarnings("unchecked")}.

\begin{oshibka}
\textbf{Частая ошибка.} Проверка \texttt{obj instanceof List<String>} не
компилируется: параметра типа в байт-коде уже нет, проверять нечего. Писать
можно \texttt{obj instanceof List<?>}~--- «какой-то список», без указания типа
элементов.
\end{oshibka}

\subsection{Подстановочные знаки и правило PECS}

Обобщения \term{инвариантны}: \texttt{List<Integer>} не является наследником
\texttt{List<Number>}, хотя \texttt{Integer} наследует \texttt{Number}. Иначе
через ссылку типа \texttt{List<Number>} в список целых можно было бы положить
дробное число. Но тогда метод, считающий сумму, пришлось бы писать отдельно для
каждого типа. Выход~--- \term{подстановочный знак} (\eng{wildcard}) \texttt{?}.

Вариантов три. Неограниченный \texttt{List<?>} принимает список любого типа, но
читать элементы позволяет только как \texttt{Object}. Ограниченный сверху
\texttt{List<? extends Number>} принимает список чисел любого конкретного типа
и разрешает читать их как \texttt{Number}, зато добавлять в него нельзя:
компилятор не знает, что там внутри~--- целые или дробные. Ограниченный снизу
\texttt{List<? super Integer>} наоборот: добавлять \texttt{Integer} безопасно.

Какой знак когда применять, подсказывает мнемоника \term{PECS}~---
\eng{Producer Extends, Consumer Super}: структура, которая отдаёт значения,
объявляется через \texttt{extends}, а та, которая их принимает,~--- через
\texttt{super} (рис.~\ref{fig:04-1}).

\begin{figure}[htbp]\centering
\begin{tikzpicture}[font=\small]
\node[boxw=40mm] (s) {\texttt{List<? extends T>}\\[1pt]
  \scriptsize производитель};
\node[boxw=26mm,right=10mm of s] (m) {\texttt{copy(...)}};
\node[boxw=36mm,right=10mm of m] (d) {\texttt{List<? super T>}\\[1pt]
  \scriptsize потребитель};
\draw[flowarrow] (s) -- (m) node[midway,above,font=\scriptsize] {читаем};
\draw[flowarrow] (m) -- (d) node[midway,above,font=\scriptsize] {пишем};
\end{tikzpicture}
\caption{Правило PECS: откуда читают и куда пишут}\label{fig:04-1}
\end{figure}

Оба ограниченных варианта показаны в листинге~\ref{lst:04-6}. Особенно
поучительна закомментированная строка в \texttt{sumNumbers}: она не
компилируется именно потому, что тип элементов списка неизвестен.

\begin{listing}[H]
\begin{minted}{java}
// Производитель: читаем числа, добавлять нельзя
public static double sumNumbers(List<? extends Number> list) {
    double sum = 0;
    for (Number n : list) sum += n.doubleValue();
    // list.add(5);  // ошибка: List<Double>? List<Integer>?
    return sum;
}

// Потребитель: пишем целые, читать можно лишь как Object
public static void addIntegers(List<? super Integer> list) {
    list.add(10);
}
\end{minted}
\caption{Подстановочные знаки в параметрах методов}
\label{lst:04-6}
\end{listing}

\section{Исключения}

Даже безупречно типизированный код сталкивается с тем, чего нельзя предвидеть
на этапе компиляции: файл удалён, сеть недоступна, пользователь ввёл буквы
вместо числа. \term{Исключение} (\eng{exception}) в Java~--- это объект,
описывающий такую ситуацию; выброшенное исключение прерывает нормальный ход
выполнения и передаёт управление тому, кто готов его обработать.

\subsection{Иерархия и две категории исключений}

Все исключения наследуют классу \texttt{Throwable}, от которого расходятся две
ветви (рис.~\ref{fig:04-2}). Ветвь \texttt{Error}~--- аварии самой виртуальной
машины: кончилась память, переполнился стек вызовов; ловить их бессмысленно.
Ветвь \texttt{Exception}~--- то, с чем приложение способно что-то сделать.

\begin{figure}[htbp]\centering
\begin{forest} hierarchy
[Throwable
  [Error
    [OutOfMemoryError]
    [StackOverflowError]
  ]
  [Exception
    [IOException
      [FileNotFoundException]
    ]
    [SQLException]
    [RuntimeException
      [NullPointerException]
      [IllegalArgumentException]
      [ArithmeticException]
    ]
  ]
]
\end{forest}
\caption{Иерархия классов исключений Java}\label{fig:04-2}
\end{figure}

Внутри ветви \texttt{Exception} проходит вторая, более важная граница.
\term{Проверяемые} (\eng{checked}) исключения~--- все наследники
\texttt{Exception}, кроме потомков \texttt{RuntimeException}; компилятор требует
их обработать или объявить. \term{Непроверяемые} (\eng{unchecked}) ничего не
требуют (табл.~\ref{tab:04-2}).

\begin{table}[htbp]
\caption{Проверяемые и непроверяемые исключения}
\label{tab:04-2}
\centering
\begin{tabular}{L{3.9cm}L{5.4cm}L{5.4cm}}
\toprule
Критерий & Проверяемые & Непроверяемые \\
\midrule
Предок & \texttt{Exception} (не \texttt{RuntimeException}) & \texttt{RuntimeException} или \texttt{Error} \\
\addlinespace
Компилятор требует обработки & да & нет \\
\addlinespace
Природа события & внешние обстоятельства & ошибка в коде \\
\addlinespace
Примеры & \texttt{IOException}, \texttt{SQLException} & \texttt{NullPointerException} \\
\bottomrule
\end{tabular}
\end{table}

Разница здесь философская, а не техническая. Проверяемое исключение~---
ожидаемая неприятность, о которой компилятор заставляет подумать заранее: файла
может не оказаться, и это не вина программиста. Непроверяемое~--- симптом
ошибки в коде: обращение по \texttt{null} нужно не обрабатывать, а устранять.

\subsection{Блоки try, catch и finally}

Опасный код помещают в блок \texttt{try}, обработку~--- в \texttt{catch}, а
завершающие действия~--- в \texttt{finally}. Блоки \texttt{catch}
просматриваются сверху вниз, и выполняется первый подходящий; поэтому
специфичные исключения должны идти перед общими~--- иначе до нижних блоков
управление не дойдёт, и компилятор сообщит об ошибке. Блок \texttt{finally}
выполняется всегда. С Java~7 несколько несвязанных типов ловятся одним блоком
через вертикальную черту (листинг~\ref{lst:04-7}).

\begin{listing}[H]
\begin{minted}{java}
try {
    int[] arr = new int[5];
    arr[10] = 1;                 // выход за границу массива
} catch (ArrayIndexOutOfBoundsException e) {
    // Сначала конкретное исключение
    System.out.println("Сообщение: " + e.getMessage());
} catch (RuntimeException e) {
    // Потом более общее - иначе код не соберётся
    System.out.println("Иная ошибка времени выполнения");
} finally {
    System.out.println("Выполнится в любом случае");
}

// Один обработчик на несколько несвязанных типов (Java 7+)
try {
    loadConfiguration();
} catch (IOException | ParseException e) {
    System.out.println("Не удалось прочитать настройки");
}
\end{minted}
\caption{Порядок блоков catch и multi-catch}
\label{lst:04-7}
\end{listing}

У любого исключения есть методы, отвечающие на вопрос, что случилось:
\texttt{getMessage} возвращает текст, \texttt{getClass} даёт точный тип, а
\texttt{printStackTrace} печатает цепочку вызовов.

\subsection{Операторы throw и throws, собственные исключения}

Два похожих слова делают разное. \texttt{throw} выбрасывает конкретный объект
исключения прямо сейчас. \texttt{throws} в заголовке метода лишь предупреждает:
метод может выбросить такое проверяемое исключение. Вызывающий код обязан либо
поймать его, либо объявить \texttt{throws} у себя и переложить ответственность
выше.

Когда стандартных классов не хватает, объявляют своё исключение: наследуют
\texttt{Exception} для проверяемого и \texttt{RuntimeException} для
непроверяемого. Конструктор обязательно вызывает конструктор предка, иначе
сообщение и первопричина потеряются. Конструктор с параметром
\texttt{Throwable cause} реализует \term{цепочку исключений} (\eng{exception
chaining}): низкоуровневая ошибка заворачивается в высокоуровневую, но не
исчезает~--- в распечатке стека она появится после слов \texttt{Caused by}
(листинг~\ref{lst:04-8}).

\begin{listing}[H]
\begin{minted}{java}
// Проверяемое исключение уровня приложения
class DataLoadException extends Exception {
    public DataLoadException(String message, Throwable cause) {
        super(message, cause);   // сохраняем первопричину
    }
}

class UserService {
    // throws предупреждает вызывающий код
    public String load(int id) throws DataLoadException {
        try {
            return database.findUser(id);
        } catch (SQLException e) {
            throw new DataLoadException(
                    "Не удалось загрузить пользователя " + id, e);
        }
    }
}
\end{minted}
\caption{Собственное исключение и цепочка причин}
\label{lst:04-8}
\end{listing}

В обработчике доступны и собственное сообщение (\texttt{getMessage}), и
сообщение первопричины (\texttt{getCause().getMessage()}).

Отдельное правило действует при переопределении: наследник не вправе расширять
список проверяемых исключений. Он может объявить то же исключение, заменить его
наследником (скажем, \texttt{FileNotFoundException} вместо
\texttt{IOException}), убрать \texttt{throws} вовсе~--- но не добавить новое.
Сотрудник, вышедший вместо коллеги, вправе доставить меньше хлопот, чем обещал
предшественник, но не больше: тот, кто с ним работает, о новых неприятностях не
предупреждён. Причина техническая: вызывающий код обращается через ссылку
базового типа и пишет \texttt{catch} по объявлению базового метода.

\subsection{Автоматическое закрытие ресурсов}

Файлы, сетевые соединения и подключения к базе данных нужно закрывать, причём
даже тогда, когда работа с ними оборвалась исключением. Написанный вручную блок
\texttt{finally} для этого громоздок: закрытие само может выбросить исключение,
поэтому внутри \texttt{finally} появляется ещё один \texttt{try}. С Java~7 эту
работу берёт на себя \term{try-with-resources}: ресурсы объявляют в круглых
скобках после \texttt{try}, и виртуальная машина закрывает их сама при любом
выходе из блока. Требование к классу ресурса одно~--- реализовать
\texttt{AutoCloseable} с методом \texttt{close} (листинг~\ref{lst:04-9}).

\begin{listing}[H]
\begin{minted}{java}
try (BufferedReader reader =
             new BufferedReader(new FileReader("data.txt"))) {
    String line;
    while ((line = reader.readLine()) != null) {
        System.out.println(line);
    }
} catch (IOException e) {
    System.out.println("Ошибка чтения: " + e.getMessage());
}
// reader закрыт автоматически - и при успехе, и при исключении
\end{minted}
\caption{Чтение файла через try-with-resources}
\label{lst:04-9}
\end{listing}

\begin{oshibka}
\textbf{Частая ошибка: «проглатывание» исключения.} Пустой блок
\texttt{catch (Exception e) \{ \}} выключает всю систему сообщений об ошибках:
программа не падает, ничего не печатает и продолжает считать по неверным
данным, а найти такое место по журналам невозможно. Ловите конкретный тип и,
если обработать его нечем, хотя бы запишите сообщение в журнал.
\end{oshibka}

\section{Отладка}

Рано или поздно случается неизбежное: программа компилируется, запускается, не
падает~--- и выдаёт неправильный ответ. Компилятор здесь уже не помощник.
Помогает \term{отладчик} (\eng{debugger})~--- инструмент, который останавливает
программу в любой точке и показывает её изнутри. Это как пульт у фильма: вместо
того чтобы пересматривать сцену в надежде разглядеть номер мелькнувшей машины,
вы жмёте паузу ровно на нужном кадре.

\subsection{Почему печать в консоль~--- плохой отладчик}

Первый инструмент, к которому тянется рука новичка,~--- вызов
\texttt{System.out.println}. Он работает, но цена растёт быстро: печатать
приходится то, что вы \emph{угадали} заранее, каждая новая гипотеза требует
перекомпиляции, отладочные строки забываются и уезжают в общий репозиторий, а в
цикле на десять тысяч итераций консоль превращается в кашу. Отладчик снимает
все эти проблемы разом: он показывает все переменные текущего кадра, проверяет
гипотезы на уже остановленной программе, живёт в настройках среды, а не в
исходниках, умеет срабатывать по условию и показывает цепочку вызовов. Печать
остаётся полезной для серверов, к которым не подключиться отладчиком: там её
роль берёт на себя журналирование.

\subsection{Точки останова}

\term{Точка останова} (\eng{breakpoint})~--- метка «останови программу здесь».
В IntelliJ~IDEA её ставят щелчком по левому полю редактора напротив строки, в
VS~Code~--- тем же щелчком или клавишей~F9; запуск в режиме отладки~---
Shift+F9 и F5 соответственно. Простая точка на строке~--- только начало:
настоящая сила в её настройках, которые открывает щелчок правой кнопкой.

Самая полезная разновидность~--- \term{условная точка останова}. В поле
\textbf{Condition} пишут обычное Java-выражение, в котором доступны все
переменные, видимые в этой строке: \texttt{count > 100 \&\& sum == 0} или
\texttt{user.getId() == 42L}. Обычная точка останова~--- консьерж, который
звонит вам про каждого входящего в подъезд; условная~--- консьерж, которому вы
сказали: «звоните, только если придёт курьер». Когда условие писать не по
чему~--- данные внешне одинаковые, а ломается 4783-я запись,~--- указывают
число пропускаемых попаданий. Есть и другие виды: точка на исключении
срабатывает в момент выброса, а логирующая печатает выражение, не останавливая
программу.

\subsection{Пошаговое выполнение и просмотр данных}

Программа остановилась~--- дальше вы двигаете её вручную. Команд всего пять, и
различаются они лишь тем, как поступают с вызовом метода на текущей строке:
выполнить его целиком, войти внутрь или выйти наружу. Горячие клавиши в
табл.~\ref{tab:04-3} даны для IntelliJ~IDEA и VS~Code.

\begin{table}[htbp]
\caption{Команды пошагового выполнения}
\label{tab:04-3}
\centering
\begin{tabular}{L{3.0cm}L{2.4cm}L{2.4cm}L{6.4cm}}
\toprule
Команда & IntelliJ IDEA & VS Code & Действие \\
\midrule
Step Over & F8 & F10 & выполнить строку целиком, внутрь вызовов не заходя \\
\addlinespace
Step Into & F7 & F11 & войти внутрь вызываемого метода \\
\addlinespace
Step Out & Shift+F8 & Shift+F11 & досчитать метод до конца и вернуться к вызывающему \\
\addlinespace
Run to Cursor & Alt+F9 & Ctrl+F10 & выполнять до строки с курсором \\
\addlinespace
Resume & F9 & F5 & продолжить до следующей точки останова \\
\bottomrule
\end{tabular}
\end{table}

Рабочая лошадка~--- \textbf{Step Over}: вы идёте по своему методу, а строку
вроде \texttt{list.sort(comparator)} проходите целиком, не проваливаясь в недра
стандартной библиотеки.

Окно \textbf{Variables} показывает все переменные, видимые в текущей точке:
параметры метода, локальные переменные и \texttt{this} со всеми полями. Именно
здесь вы отвечаете на главный вопрос отладки: какие данные тут на самом деле? В
окно \textbf{Watches} добавляют выражения, за которыми хочется следить
постоянно; они пересчитываются на каждом шаге. А \textbf{Evaluate Expression}
(Alt+F8 в IDEA, консоль отладки в VS~Code) вычисляет любое выражение разово, в
контексте остановленной программы, включая вызовы методов. Значение переменной
можно и изменить, выделив её в окне Variables и нажав~F2: так проверяют
гипотезу без перекомпиляции.

\begin{oshibka}
\textbf{Осторожно с побочными эффектами.} Вычисление в Evaluate Expression
выполняется по-настоящему: если вы напишете \texttt{list.remove(0)}, элемент
действительно удалится и программа пойдёт другим путём. И помните: изменение
переменной в отладчике не исправляет код~--- это разведка, после которой всё
равно нужно править исходники.
\end{oshibka}

\subsection{Стек вызовов и точки останова на исключении}

Каждый вызов метода кладёт на стек \term{кадр} (\eng{frame}) с параметрами и
локальными переменными этого вызова; завершение метода кадр снимает. Стек
вызовов~--- стопка тарелок: каждый вызов кладёт тарелку сверху, возврат снимает
её; достать из середины нельзя, а посмотреть на любую~--- можно. Окно
\textbf{Frames} показывает всю цепочку от текущего метода до \texttt{main}, а
щелчок по кадру переключает окно Variables на переменные \emph{того} вызова.
Приём незаменим, когда метод получил некорректный аргумент: виноват не он, а
тот, кто его вызвал.

Распечатка стека в консоли отвечает на вопрос «где упало», но молчит про «с
какими данными»: к моменту печати кадры уже сняты. Отладчик умеет остановиться
в момент выброса, пока состояние ещё живо: в IDEA для этого в списке точек
останова добавляют \textbf{Java Exception Breakpoint} и вписывают класс
исключения. Режим \textbf{Caught} останавливает даже на тех исключениях,
которые кто-то ловит,~--- именно он находит «проглоченные» исключения из
пустого \texttt{catch}, а их иначе не найти вовсе.

\subsection{Разбор на примере}

Возьмём программу с настоящей ошибкой (листинг~\ref{lst:04-10}). Она считает
среднее по списку строк, пропуская те, что не разбираются в число. Ожидаем
\texttt{20.0}~--- среднее от 10, 20 и 30,~--- а получаем \texttt{15.0}:
программа не упала, исключение поймано, ответ неверный. Идеальный случай для
отладчика.

\begin{listing}[H]
\begin{minted}{java}
public static double average(List<? extends CharSequence> rows) {
    double sum = 0;
    int count = 0;
    for (CharSequence row : rows) {
        try {
            sum += Double.parseDouble(row.toString().trim());
        } catch (NumberFormatException e) {
            // Некорректную строку просто пропускаем
        }
        count++;
    }
    return sum / count;
}

List<String> rows = List.of("10", "20", "abc", "30");
System.out.println("Среднее: " + average(rows));   // 15.0
\end{minted}
\caption{Программа, которая считает неправильно}
\label{lst:04-10}
\end{listing}

Ставим точку останова на строке \texttt{return sum / count;} и запускаем
отладку. Окно Variables показывает \texttt{sum = 60.0} и \texttt{count = 4}:
сумма верна (10 + 20 + 30), а счётчик равен четырём, хотя чисел было три.
Проверяем гипотезу, не трогая код: в Evaluate Expression вычисляем
\texttt{sum / (count - 1)} и получаем \texttt{20.0}.

Теперь выясним, что происходит на строке \texttt{"abc"}. Ставим точку останова
на исключении \texttt{NumberFormatException} в режиме Caught и перезапускаем
отладку. Программа замирает не в нашем коде, а внутри стандартной библиотеки,
куда её увёл \texttt{Double.parseDouble}; в окне Frames щёлкаем по первому
своему кадру~--- \texttt{average}~--- и видим \texttt{row} равным
\texttt{"abc"}, \texttt{sum} равной \texttt{30.0}, \texttt{count} равным~2.
Нажимаем Step Out, пока в редакторе не окажется наш метод: исключение как раз
дойдёт до блока \texttt{catch}. Дальше Step Over~--- блок \texttt{catch} пуст, и
следующей выполняется строка \texttt{count++}. Счётчик становится равным трём,
хотя строка \texttt{"abc"} в сумму не попала. Ошибка найдена: инкремент стоит
\emph{вне} блока \texttt{try}, поэтому считаются все строки.

Остаётся подтвердить диагноз: доходим до строки возврата, ставим \texttt{count}
равным трём клавишей~F2, продолжаем выполнение~--- и программа печатает
\texttt{20.0}. Исправление очевидно: строку \texttt{count++} надо перенести
внутрь блока \texttt{try}, сразу после успешного разбора. Заодно стоит закрыть
две проблемы, которые отладка вывела на свет: пустой \texttt{catch} должен хотя
бы печатать сообщение о пропущенной строке, а перед делением нужна проверка
\texttt{if (count == 0)} с выбросом \texttt{IllegalArgumentException}~--- иначе
при полном отсутствии чисел получится деление \texttt{0.0} на~0, а это для
\texttt{double} не исключение, а тихое \texttt{NaN}, отравляющее все дальнейшие
расчёты.

\subsection{Практические правила отладки}

Сначала воспроизведите ошибку: пока вы не умеете вызывать её по требованию,
отлаживать нечего. Формулируйте гипотезу до нажатия клавиши~--- «проверяю, что
\texttt{count} равен четырём» есть отладка, «пощёлкаю и посмотрю» есть трата
времени. Ставьте точку останова не там, где падает, а там, где данные ещё
правильные: ошибка почти всегда рождается раньше, чем проявляется. И не
отлаживайте вхолостую: условие или счётчик попаданий дешевле сотни нажатий
Resume.

\praktikum{}

Понадобится JDK~17 или новее и среда разработки с отладчиком. Каждое задание
оформляется отдельным классом; в отчёт входят исходный код, результат запуска и
письменные ответы на поставленные вопросы. Фрагменты в листингах теории
иллюстративны: многоточия и вызовы вроде \texttt{database.findUser} заменяйте
собственными заглушками.

\subsection*{Задание 1. Вложенные классы}

Реализуйте класс \texttt{Library} с полями \texttt{name} и \texttt{capacity}, а
внутри него~--- нестатический внутренний класс \texttt{Book} с полями
\texttt{title}, \texttt{author}, \texttt{year} и методом \texttt{getInfo},
печатающим строку вида «Книга: Война и мир автора Толстой~Л.~Н. (1869) в
библиотеке Городская библиотека»; \texttt{Book} должен обращаться к полю
\texttt{name} внешнего класса напрямую. Создайте объект через
\texttt{lib.new Book(...)} и вызовите \texttt{getInfo}. После этого добавьте к
классу \texttt{Book} слово \texttt{static} и объясните письменно, какая ошибка
компиляции возникла и почему потерялся доступ к полю \texttt{name}.

Вторая часть~--- строитель. Реализуйте класс \texttt{Computer} с полями
\texttt{final}: \texttt{cpu}, \texttt{ram}, \texttt{storage} (по умолчанию 256),
\texttt{gpu} (по умолчанию \texttt{null}) и \texttt{ssd} (по умолчанию
\texttt{true}). Конструктор сделайте закрытым и принимающим строитель, а
статический вложенный класс \texttt{Builder} снабдите конструктором с двумя
обязательными параметрами, методами настройки остальных полей, возвращающими
\texttt{this}, и методом \texttt{build}. Переопределите \texttt{toString} в
формате \texttt{Computer\{CPU='...', RAM=...ГБ, накопитель=...ГБ,
видеокарта=...\}}: без него \texttt{System.out.println} напечатает служебную
строку вида \texttt{Computer@1b6d3586}. Соберите игровой компьютер (Core~i9,
32~ГБ, 2000~ГБ, RTX~4080) и офисный (Core~i5, 16~ГБ) и напечатайте оба.

Третья часть~--- вложенный интерфейс. Объявите внутри \texttt{Library}
интерфейс \texttt{BookFilter} с методом \texttt{boolean accept(Book b)} и метод
\texttt{printFiltered}, печатающий только те книги, для которых фильтр вернул
\texttt{true}. Фильтр реализуйте отдельным классом (скажем, «книги новее
1900~года»), обратившись к контракту как \texttt{Library.BookFilter}.

\subsection*{Задание 2. Обобщённый стек и подстановочные знаки}

Реализуйте обобщённый стек \texttt{Stack<T>} на основе массива
\texttt{Object[]}: начальная ёмкость~10, при переполнении массив расширяется
вдвое. Методы \texttt{push}, \texttt{pop} и \texttt{peek} (два последних
бросают \texttt{EmptyStackException} из пакета \texttt{java.util} на пустом
стеке), а также \texttt{isEmpty}, \texttt{size} и \texttt{toString}.
Продемонстрируйте работу на
\texttt{Stack<String>} и \texttt{Stack<Integer>}.

Вторая часть~--- класс \texttt{NumberUtils} со статическими методами
\texttt{sum(List<? extends Number>)} и \texttt{average(List<? extends Number>)},
где второй бросает \texttt{IllegalArgumentException} на пустом списке;
\texttt{<T extends Comparable<T>> T findMax(List<T>)} и парный ему
\texttt{findMin}; \texttt{fillWithIntegers(List<? super Integer>, int n)},
добавляющий числа от~1 до~$n$. Проверьте методы на списках чисел и строк и
объясните письменно, почему \texttt{List<Integer>} нельзя присвоить переменной
типа \texttt{List<Number>} и как подстановочный знак решает эту задачу.

\subsection*{Задание 3. Исключения банковского приложения}

Постройте систему исключений для банка. Базовое проверяемое исключение
\texttt{BankException} наследует \texttt{Exception} и хранит поле
\texttt{operationId}; наследник \texttt{InsufficientFundsException} принимает
требуемую и доступную суммы и формирует из них сообщение; непроверяемое
\texttt{InvalidAmountException} наследует \texttt{RuntimeException}.

Класс \texttt{BankAccount} хранит \texttt{accountId} и \texttt{balance}. Его
конструктор бросает \texttt{InvalidAmountException} при отрицательном начальном
балансе; метод \texttt{deposit} бросает его же при сумме, не большей нуля;
метод \texttt{withdraw} объявляет \texttt{throws InsufficientFundsException} и
бросает это исключение, когда денег не хватает. Напишите метод \texttt{main},
проверяющий все четыре ситуации, и объясните в отчёте, почему одни исключения
сделаны проверяемыми, а другие~--- нет.

Третья часть~--- ресурсы и цепочка причин. Класс \texttt{TransactionLog}
реализует \texttt{AutoCloseable}: конструктор печатает «журнал открыт», метод
\texttt{write}~--- номер и текст записи, метод \texttt{close}~--- «журнал
закрыт». Файл здесь не нужен: важно увидеть,
что \texttt{close} вызывается сам и при обычном выходе, и при исключении внутри
блока. Затем пусть \texttt{write} на каждой пятой записи бросает
\texttt{IOException} (пакет \texttt{java.io}), а вызывающий код ловит его и
заворачивает в \texttt{BankException} с сохранением причины; напечатайте и
\texttt{getMessage}, и \texttt{getCause().getMessage()}.

\subsection*{Задание 4. Поиск ошибки отладчиком}

Наберите программу из листинга~\ref{lst:04-10} и повторите разбор
самостоятельно, не подглядывая в текст главы. Убедитесь, что печатается
\texttt{15.0}; поставьте точку останова на строке возврата и запишите значения
\texttt{sum} и \texttt{count}. Пройдите цикл командой Step Over итерацию за
итерацией, фиксируя \texttt{row}, \texttt{sum} и \texttt{count}; на строке с
\texttt{Double.parseDouble} один раз примените Step Into, а затем вернитесь в
свой код командой Step Out. Поставьте на строке \texttt{count++}
условную точку останова с условием \texttt{row.toString().equals("abc")} и
проверьте, что отладка останавливается сразу на нужной итерации. Наконец,
вычислите в Evaluate Expression выражение \texttt{sum / (count - 1)}.

Ответьте письменно: после какой именно строки тела цикла значения \texttt{sum}
и \texttt{count} перестают соответствовать ожиданиям; почему \texttt{sum} верна,
а \texttt{count}~--- нет. Исправьте метод и приложите результат повторного
запуска.

\subsection*{Результаты занятия}

По итогам занятия у вас должны быть готовы: файлы \texttt{.java} со всеми
четырьмя заданиями и результатами их запуска; объяснение ошибки компиляции при
добавлении \texttt{static} к внутреннему классу; обоснование выбора
подстановочных знаков в \texttt{NumberUtils}; разбор того, какие банковские
исключения сделаны проверяемыми и почему; протокол отладки из задания~4 с
исправленным кодом.

\kontrolvoprosy

\begin{enumerate}
\item Чем нестатический внутренний класс отличается от статического вложенного?
      Почему для создания объекта первого нужен объект внешнего класса?
\item Почему строитель реализуют статическим вложенным классом, а итератор
      коллекции~--- внутренним?
\item Какие неявные модификаторы получает интерфейс, объявленный внутри другого
      интерфейса?
\item Что такое стирание типов и какие операции оно запрещает по отношению к
      параметру типа?
\item Сформулируйте правило PECS. Почему в список типа
      \texttt{List<? extends Number>} нельзя ничего добавить?
\item Чем проверяемые исключения отличаются от непроверяемых и какая философия
      стоит за этим делением?
\item Почему специфичные блоки \texttt{catch} должны идти раньше общих?
\item В чём разница между \texttt{throw} и \texttt{throws}? Какие изменения
      списка \texttt{throws} разрешены при переопределении метода?
\item Для чего нужна цепочка исключений и что теряется, если не передать
      причину в конструктор предка?
\item Чем условная точка останова удобнее обычной и зачем нужна точка останова
      на исключении?
\end{enumerate}

\testzadaniya

\begin{enumerate}

\item Как создать объект нестатического внутреннего класса \texttt{Inner},
      объявленного внутри класса \texttt{Outer}?
  \begin{enumerate}
    \item[а)] \texttt{Inner i = new Inner();}
    \item[б)] \texttt{Outer.Inner i = new Outer.Inner();}
    \item[в)] \texttt{Outer o = new Outer(); Outer.Inner i = o.new Inner();}
    \item[г)] \texttt{Inner i = Outer.new Inner();}
  \end{enumerate}

\item К каким членам внешнего класса имеет доступ статический вложенный класс?
  \begin{enumerate}
    \item[а)] ко всем, включая нестатические \texttt{private}
    \item[б)] только к статическим членам внешнего класса
    \item[в)] только к членам с модификатором \texttt{public}
    \item[г)] ни к каким~--- он полностью изолирован
  \end{enumerate}

\item Какую основную задачу решают обобщения?
  \begin{enumerate}
    \item[а)] ускоряют выполнение программы
    \item[б)] дают типобезопасность на этапе компиляции и избавляют от
              приведения типов
    \item[в)] позволяют хранить примитивы в коллекциях без упаковки
    \item[г)] автоматически преобразуют объекты в текстовый формат
  \end{enumerate}

\item Какое действие невозможно из-за стирания типов?
  \begin{enumerate}
    \item[а)] объявить переменную типа \texttt{T}
    \item[б)] передать объект типа \texttt{T} в метод
    \item[в)] объявить поле типа \texttt{List<T>}
    \item[г)] выполнить \texttt{new T()} или \texttt{new T[10]}
  \end{enumerate}

\item Правило PECS подсказывает: если из коллекции только читают, параметр
      объявляют как\ldots
  \begin{enumerate}
    \item[а)] \texttt{List<? super T>}
    \item[б)] \texttt{List<T>}
    \item[в)] \texttt{List<? extends T>}
    \item[г)] \texttt{List<Object>}
  \end{enumerate}

\item Чем проверяемые исключения отличаются от непроверяемых?
  \begin{enumerate}
    \item[а)] непроверяемые нельзя перехватить в блоке \texttt{catch}
    \item[б)] проверяемые возникают только во время выполнения
    \item[в)] непроверяемые наследуют \texttt{Exception} напрямую
    \item[г)] проверяемые обязаны быть обработаны или объявлены в
              \texttt{throws}
  \end{enumerate}

\item Блок \texttt{try} печатает \texttt{A}, затем выполняет деление на ноль и
      печатает \texttt{B}; блок \texttt{catch} печатает \texttt{C}, блок
      \texttt{finally}~--- \texttt{D}. Что окажется в консоли?
  \begin{enumerate}
    \item[а)] \texttt{ABCD}
    \item[б)] \texttt{AD}
    \item[в)] \texttt{ACD}
    \item[г)] \texttt{AB}
  \end{enumerate}

\item Что называют цепочкой исключений?
  \begin{enumerate}
    \item[а)] несколько блоков \texttt{try-catch}, идущих подряд
    \item[б)] несколько операторов \texttt{throw} в одном методе
    \item[в)] перехват нескольких типов одним блоком \texttt{catch}
    \item[г)] обёртывание исходного исключения в новое с передачей причины
  \end{enumerate}

\item Чем условная точка останова отличается от обычной?
  \begin{enumerate}
    \item[а)] она срабатывает лишь тогда, когда истинно заданное выражение, а
              обычная~--- при каждом проходе строки
    \item[б)] она работает только с числовыми переменными
    \item[в)] она доступна только в одной из сред разработки
    \item[г)] её нельзя снять после установки
  \end{enumerate}

\item При остановке на строке \texttt{return sum / count;} окно Variables
      показало \texttt{sum = 60.0} и \texttt{count = 4}, хотя корректных чисел
      в списке было три. О чём это говорит?
  \begin{enumerate}
    \item[а)] переменная \texttt{sum} вычислена неверно
    \item[б)] счётчик увеличивается для каждой строки, включая неразобранную
    \item[в)] список содержит повторяющиеся элементы
    \item[г)] отладчик не видит исключений внутри библиотеки
  \end{enumerate}

\end{enumerate}
